\documentclass[11pt]{article}
\usepackage[margin=2.5cm]{geometry}
\usepackage{amsmath,amssymb,amsthm,mathtools}
\usepackage{bm}
\usepackage{booktabs}
\usepackage[round]{natbib}
\usepackage{hyperref}

\newcommand{\R}{\mathbb{R}}
\newcommand{\Z}{\mathbb{Z}}
\newcommand{\E}{\mathbb{E}}
\newcommand{\Prob}{\mathbb{P}}
\newcommand{\neff}{n_{\mathrm{eff}}}
\newcommand{\dip}{\mathrm{dip}}
\renewcommand{\thesection}{S\arabic{section}}
\renewcommand{\thetable}{S\arabic{table}}
\renewcommand{\theequation}{S\arabic{equation}}

\title{Supplementary material for\\
``Effective-sample-size calibration of a modality test for cluster
homogeneity and the number of classes in spatially autocorrelated data''}
\author{Francisco Parra}
\date{}

\begin{document}
\maketitle

\noindent
This supplement collects three things that would weigh down the main text: the
full discussion of what Proposition~1 does and does not establish (Section~S1);
a complete specification of every criterion in the benchmark, sufficient to
reproduce the comparison (Section~S2); and the implementation of ESS-Dip in
enough detail to re-code it from the description alone (Section~S3). Section,
equation and table numbers without an ``S'' refer to the main text. The
released code is the executable reference and fixes every seed.

\section{What Proposition 1 does and does not establish}\label{sec:supp-prop}

\subsection*{Proof of Proposition 1}

For convenience we restate what is proved. Under (C1)--(C3), with $C$
increasing to $\Z^{2}$ along a sequence of rectangles and the test calibrated
at $\neff=n_C/a_I$, the rejection probability of the per-node test tends to
zero, and hence to at most $\alpha$; consequently $\Prob(\hat K=1)\to1$ when
every node visited satisfies (C1).

\begin{proof}
The argument reduces the dependent problem to the independent one; the
spatial dependence enters only through \textnormal{(C3)}, and the properties
of the dip that are used involve no assumption on the dependence structure of
the data. First, since $\Phi$ is unimodal, the deterministic inequality
$\hat D_C=\dip(\hat F_C)\le\sup_x\lvert\hat F_C(x)-\Phi(x)\rvert$
\citep{HartiganHartigan1985} together with condition (C3) shows that
$\sqrt{\neff}\,\hat D_C$ is tight: the dip vanishes at least at the
$\neff^{-1/2}$ rate. Second, at a strictly unimodal law the rate is strictly
faster. Fluctuations of size $\epsilon$ about a distribution function with
strictly curved branches can be absorbed into the unimodal class to within
$o(\epsilon)$, an absorption that uses only the uniform continuity of the
limiting process. For independent samples this is known in two forms:
\citet[Theorem~5]{HartiganHartigan1985} prove $\sqrt m\,D_m\to0$ in
probability for unimodal laws possessing a nonzero derivative of some order
$k\ge2$ at the mode and an exponentially decreasing density, a class they note
contains the normal; and \citet[Theorem~2]{ChengHall1998} sharpen the rate to
$D_m=O_p(m^{-3/5})=o_p(m^{-1/2})$ under their condition~(4.2), which the
Gaussian also satisfies. Under (C3) the same absorption gives
$\sqrt{\neff}\,\hat D_C\to0$ in probability. Third, the calibration
quantile inherits $q_{1-\alpha}(m)=c_\alpha m^{-1/2}(1+o(1))$, with
$c_\alpha>0$, from the uniform law, for which $\sqrt m\,D_m$ converges to the
nondegenerate dip of a Brownian bridge
\citep[Theorem~3]{HartiganHartigan1985}. Hence
$\Prob\bigl(\hat D_C>q_{1-\alpha}(\lfloor\neff\rceil)\bigr)\to0$, which
gives the bound of the proposition, and indeed a limiting rejection probability of zero;
the same conclusion holds for any projected law in the class of
\citet[Theorem~5]{HartiganHartigan1985}. The bound $\alpha$ itself is
approached only in the flat-topped boundary case, where linear stretches of the
distribution function leave no curvature to absorb the fluctuations; that the
uniform is least favourable among \emph{all} unimodal laws, and not only
within that class, is stated asymptotically by
\citet{HartiganHartigan1985} but proved by them only for laws with
exponentially decreasing densities, so we invoke the bound in the form their
theorems establish. Under the single-class
hypothesis every node visited by the recursion satisfies \textnormal{(C1)};
the recursion visits finitely many nodes, so a union bound over their
(vanishing) rejection probabilities gives $\Prob(\hat K=1)\to1$.
\end{proof}

\subsection*{Scope}

Proposition~1 concerns a stylised version of the test: the projection direction
is held fixed, the detrending step is switched off, and the decorrelation area
is the indicator integral range $a_I$ rather than the deployed $\ell^{2}$
proxy. It is a consistency check on the calibration principle, not a statement
about the procedure as applied. Four gaps separate the two, and we set them out
in turn.

\paragraph{(i) Condition (C3) carries the analytic weight.}
Condition (C3) is the empirical-process counterpart, at the indicator integral
range, of the effective-sample-size heuristic. It therefore assumes precisely
what one would most want proved: that ``plugging an effective sample size into
the independent reference distribution'' is legitimate for a distributional
statistic such as the dip, and that the appropriate decorrelation scale for the
dip is the indicator range $a_I$ rather than the field integral range that
governs the sample mean (Section~3.7). Functional central limit theorems of
this form are established for the empirical processes of strongly mixing and of
associated sequences and fields, and condition (C1) places $\bm\varepsilon$ in
the associated class through \citet{Pitt1982} and
\citet{BulinskiShashkin2007}; a complete verification for the lattice empirical
process under our exact conditions is beyond the present scope, so (C3) is
stated as a hypothesis and not derived. A reader who doubts (C3) should regard
Proposition~1 as showing that the rest of the argument is sound \emph{given}
the effective-sample-size premise, and nothing more.

\paragraph{(ii) Why the level is conservative rather than exact.}
The critical value is calibrated on the least-favourable uniform null, whereas
under (C1) the projected data are Gaussian, whose rescaled dip vanishes
asymptotically; the rejection probability therefore tends to zero rather than
to $\alpha$. \citet{ChengHall1998} note the same degeneracy for
uniform-calibrated dip and excess-mass tests. This is consistent with the zero
empirical false-split rate of Figure~1a. Attaining level $\alpha$ would require
calibrating on the model's own null, which is what the recall variant
ESS-Dip-R does, at the price of the distribution-free property.

\paragraph{(iii) The deployed default is assessed empirically, not covered by
the proposition.}
The proposition takes the decorrelation area to be $a_I$. The deployed default
estimates it by the $\ell^{2}$ proxy of Section~3.5, which in our experiments
lies below $a_I$ ($\ell^{2}\approx0.3$--$0.56\,a_I$) and therefore
\emph{overstates} the effective sample size, pushing the test towards
anti-conservatism. It is the slack of point~(ii) that absorbs this bias. The
net level of the deployed pipeline is an empirical matter, addressed by
Figure~1, by the range-perturbation experiment of Section~4.8, and by the
single-class windows of Section~5.3. The same accounting explains why
ESS-Dip-R, which removes the uniform slack, is not covered by the proposition
at all and is assessed on the model's null only.

\paragraph{(iv) Detrending and the projection direction.}
Both are idealisations, but they are not equally consequential. The projection
direction is data-dependent in practice, (C2) being the fixed-direction
idealisation; Figure~1a shows no inflation of the empirical level when the
$2$-means direction is used instead, but this is evidence, not proof.

Switching off the detrending step, by contrast, costs no generality
asymptotically under (C1). There is no deterministic trend under the null, and
the adaptive rule of Section~3.6 is asymptotically inactive. The quadratic
design spans a fixed number $d$ of low-frequency directions ($d=6$ for degree
two in two coordinates), so the sum of squares it can capture is at most $d$
times the largest eigenvalue of the residual's covariance matrix. By the
Gershgorin bound that eigenvalue is at most
$\sum_{\bm h}\lvert C(\bm h)\rvert=\sigma^{2}a$, finite precisely because (C1)
makes the correlogram absolutely summable, while the total sum of squares grows
in proportion to $n_C\sigma^{2}$. Hence
\begin{equation}\label{eq:r2rate}
  \E\,R^{2}_{\mathrm{poly}}\;\le\;\frac{d\,\sigma^{2}a}{n_C\sigma^{2}}
  \;=\;\frac{d}{\neff},
  \qquad\text{so}\qquad
  \Prob\bigl(R^{2}_{\mathrm{poly}}>\tau\bigr)\;\le\;\frac{d}{\tau\,\neff}
  \;\longrightarrow\;0
\end{equation}
by Markov's inequality: the detrending trigger switches off at the same rate at
which the calibration becomes valid. Under long-range dependence, where
summability fails, this argument breaks down, as does the calibration itself.

Finally, non-Gaussian within-class laws fall outside (C1) altogether. There the
distribution-free uniform null of ESS-Dip is the safer reference, and ESS-Dip-R
may reject a homogeneous but skewed or heavy-tailed region.

\section{Specification of the benchmarked criteria}\label{sec:supp-methods}

All sixteen criteria search over $K=1,\dots,K_{\max}$ with $K_{\max}=8$ on the
factorial study and $K_{\max}=20$ on full real scenes. Unless stated otherwise
the partitions scored are produced by $k$-means. Every stochastic element is
seeded; the released code reproduces each number in the paper.

\paragraph{Partitions.}
For the four standard indices and the gap statistic, partitions come from
MiniBatch $k$-means \citep{Sculley2010} with three initialisations
(\texttt{n\_init}${}=3$) and batch size $256$, the best of the three by the
$k$-means objective being retained. Within ESS-Dip and its relatives the
binary split at each node is ordinary $k$-means with five initialisations
(\texttt{n\_init}${}=5$), again keeping the best by objective.

\paragraph{Gap statistic \citep{Tibshirani2001}.}
$\log W_K$ is computed on all pixels for $K=1,\dots,8$. Eight reference
datasets of the same size are drawn uniformly and independently over the
per-band bounding box of the data, that is, coordinate $b$ is drawn uniformly
on $[\min_s z_b(s),\max_s z_b(s)]$. With
$\mathrm{gap}(K)=\overline{\log W^{\ast}_K}-\log W_K$ and
$s_K=\mathrm{sd}(\log W^{\ast}_K)\sqrt{1+1/8}$, the estimate is the smallest
$K$ with $\mathrm{gap}(K)\ge\mathrm{gap}(K+1)-s_{K+1}$, the one-standard-error
rule of the original paper.

\paragraph{Average silhouette \citep{Rousseeuw1987}, Calinski--Harabasz
\citep{CalinskiHarabasz1974}, Davies--Bouldin \citep{DaviesBouldin1979}.}
Each index is evaluated for $K=2,\dots,8$; none can return $K=1$ by
construction, an asymmetry discussed in Section~4.3. Because the silhouette
requires all pairwise distances, all three are computed on a common uniform
subsample of $\min(1500,N)$ pixels drawn without replacement, which keeps the
$O(m^{2})$ cost bounded and applies the same sample to the three indices; the
subsample is redrawn per scene, not per $K$. The estimate maximises the index
($-$Davies--Bouldin, since lower is better there).

\paragraph{Gap with an autocorrelation-matched reference (Section~4.2).}
Identical to the gap statistic except for the reference datasets, which are
drawn as per-band Gaussian random fields whose $1/e$ range matches the
tile-median estimate $\ell$ of the scene (Section~3.5) and whose per-band
standard deviation matches the \emph{estimated} standard deviation of the
corresponding band, both quantities being computed from the data at hand, as
any deployable version of the method must. The generator is the Gaussian-kernel
construction of Section~4.1 with bandwidth $\sigma=\ell/2$, for which the
correlogram $\exp(-h^{2}/4\sigma^{2})$ crosses $1/e$ at $h=2\sigma=\ell$.

\paragraph{Calibrated validity index \citep{HennigLin2015}, two variants.}
The calibrated index is
\begin{equation}\label{eq:hl}
  V(K)=\frac{\mathrm{ASW}(K)-\overline{\mathrm{ASW}}_{\mathrm{null}}(K)}
             {\mathrm{sd}_{\mathrm{null}}(K)},
\end{equation}
where $\mathrm{ASW}$ is the average silhouette width and the null moments are
estimated from $B=20$ null datasets. The estimate is
$\hat K=\arg\max_K V(K)$ when $\max_K V(K)>1.96$ and $\hat K=1$ otherwise; the
threshold $1.96$ is the two-sided standard-normal quantile at level $0.05$, and
Section~4.5 reports the sensitivity of the comparison to it. As for the
standard indices, each dataset is reduced to a uniform subsample of $1200$
pixels before the silhouette is computed, so that $20$ null replicates remain
tractable at every $K$. The two variants differ only in the null: the
\emph{i.i.d.} variant draws each band independently from a Gaussian matched to
the band's mean and variance; the \emph{spatial} variant draws phase-randomised
surrogates, which preserve each band's periodogram, hence its variogram, while
destroying any phase structure.

\paragraph{Dip-means at the nominal size \citep{Kalogeratos2012}.}
The recursion and the split direction are those of Algorithm~1, but the dip is
referred to the uniform null at the nominal size $n_C$ rather than at
$\neff$, and neither the local range nor the trend removal is applied. This
isolates the contribution of the spatial machinery.

\paragraph{ESS-Dip, ESS-Dip-R and the four ablations.}
Section~S3 specifies these. The ablations disable one component each: the
\emph{global range} variant estimates $\ell$ over the whole image instead of as
a tile median; \emph{no trend removal} sets $\tau=\infty$; \emph{forced trend
removal} sets $\tau=0$; and \emph{declustering} replaces the calibration by
subsampling the scene to one pixel per decorrelation patch (a lattice of
spacing $\ell$, averaged over six offsets) and applying the uncalibrated dip
test to the thinned sample, with a node tested only if it retains at least $25$
thinned pixels and each child at least $12$.

\paragraph{SGS-Dip, with and without a multiple-testing correction.}
The per-node quantile is obtained from $B=40$ unconditional Gaussian
simulations on the support of the candidate cluster, generated by the spectral
method with the variogram fitted to that node's projected values. The corrected
variant tests each node at level $\alpha/K_{\max}$, with the number of
simulations raised to resolve the smaller quantile. Prespecifying the number of
tests is possible here only because the recursion is capped: a binary recursion
that returns at most $K_{\max}$ leaves performs at most $K_{\max}-1$ splits and
therefore visits at most $2K_{\max}-1$ nodes, so $\alpha/K_{\max}$ is a valid,
if conservative, Bonferroni level for the family of tests actually run. Without
the cap the number of tests would not be bounded in advance and the correction
would have no prespecified denominator.

\section{Implementation of ESS-Dip}\label{sec:supp-impl}

\paragraph{Inputs.}
The method takes a feature image $\bm z$ on a lattice $D$ with $p$ bands. No
preprocessing is part of the method. In the real-data experiments bands are
standardised to zero mean and unit variance and the hyperspectral cubes are
reduced by principal components to $p=10$; Section~5.3 reports the
insensitivity of the results to truncations between five and thirty components.

\paragraph{Constants.}
Every constant is fixed once across all experiments; none is tuned per scene.

\begin{center}
\small
\begin{tabular}{llp{6.4cm}}
\toprule
Symbol & Value & Role\\
\midrule
$\alpha$ & $0.05$ & level of the per-node test, Eq.~(6)\\
$\tau$   & $0.25$ & threshold on $R^{2}_{\mathrm{poly}}$ that triggers
                    detrending, Eq.~(11)\\
--       & $2$    & degree of the polynomial drift surface, Eq.~(10)\\
$t$      & $24$   & tile side in pixels, reduced to
                    $\min(24,\lfloor H/2\rfloor,\lfloor W/2\rfloor)$ on domains
                    too small for it\\
$\nu$    & $12$   & smallest $\neff$ at which a node is tested; below it the
                    node becomes a leaf untested\\
$B$      & $5$    & number of independent runs of the recursion whose median
                    is reported (distinct from the $B$ simulations per node of
                    the Gaussian-simulation reference)\\
--       & $300$  & Monte Carlo draws behind each null quantile\\
--       & $[12,2000]$ & interval to which $\lfloor\neff\rceil$ is clamped
                    before the quantile is looked up\\
$K_{\max}$ & $8$ / $20$ & cap on the number of leaves (factorial study / full
                    scenes)\\
\bottomrule
\end{tabular}
\end{center}

\paragraph{Range estimation.}
The decorrelation area $a=\ell^{2}$ is estimated once per scene, not per node,
and on the detrended field when the rule of Section~3.6 fires. The lattice is
cut into disjoint $t\times t$ tiles, discarding any remainder at the right and
bottom edges. Within tile $j$ the first principal component of the $p$ bands is
taken, its two-dimensional sample autocorrelation is computed as the inverse
Fourier transform of the periodogram, and the profiles along the two axes are
averaged; $\ell_j$ is the first lag at which the averaged profile falls below
$1/e$, and $\ell$ is the median of $\{\ell_j\}$. If a tile's profile never
falls below $1/e$ within the tile, that tile contributes $\ell_j=t$; if no
tile is available because the domain is smaller than $t$ in both directions,
the estimate falls back to the whole-domain range.

\paragraph{Null quantiles.}
For a node with effective size $\neff$, let
$m=\min\{\max\{\lfloor\neff\rceil,12\},2000\}$. The critical value
$q_{1-\alpha}(m)$ is the empirical $(1-\alpha)$ quantile of $300$ dip
statistics computed from $m$ pseudo-random draws of the reference law:
$\mathrm{U}(0,1)$ for ESS-Dip and $\mathrm{N}(0,1)$ for ESS-Dip-R. Quantiles
are memoised on the pair (reference law, $m$), so each distinct size is
simulated once per run however many nodes request it; this is what keeps the
calibration cost negligible relative to the $k$-means and dip evaluations.

\paragraph{Recursion.}
The stack is initialised with the whole scene. A node $C$ is popped; if
$\neff=n_C/a<\nu$ it becomes a leaf without being tested. Otherwise $2$-means
is run on the feature vectors of $C$, the split direction is the normalised
difference of the two cluster means, and the dip is computed on the projected
values by the algorithm of \citet{HartiganHartigan1985}. If the dip exceeds
$q_{1-\alpha}(m)$ \emph{and} both children retain at least $a$ pixels, both
children are pushed; otherwise $C$ becomes a leaf. The recursion stops early if
the number of leaves plus the stack depth would reach $K_{\max}$, in which case
$K_{\max}$ is returned. The estimate is the number of leaves.

\paragraph{Ensemble.}
The recursion is run $B=5$ times with independent seeds and the median of the
five values of $\hat K$ is reported. The randomness being averaged over is that
of the $k$-means initialisations and of the Monte Carlo quantiles; note that
the choice among initialisations \emph{within} a node is already made by the
$k$-means objective (\texttt{n\_init}${}=5$), so the ensemble absorbs residual
variation in the path the recursion takes, not variation in the quality of any
single split.

\paragraph{Cost.}
Each node costs one $k$-means on its pixels, one projection and one dip
evaluation, the last being $O(n_C\log n_C)$ through the sort. Range estimation
is one FFT per tile, $O(N\log t)$ over the scene. The null quantiles are
amortised by memoisation. Total cost is therefore near-linear in $N$ up to the
depth of the recursion, which is bounded by $K_{\max}$; Section~4.6 reports the
measured scaling.

\bibliographystyle{plainnat}
\bibliography{refs}
\end{document}
